\documentclass[10pt,a4paper]{amsart}
\usepackage[margin=19mm]{geometry}
\usepackage{amsmath,amssymb,mathtools}
\usepackage[scr=rsfs,cal=euler]{mathalfa}
\usepackage{xfrac}
\usepackage[unicode,hidelinks,bookmarks=false]{hyperref}
\newcommand{\ie}{i.e.\ }
\newcommand{\cf}{cf.\ }
\newcommand{\resp}{respectively\ }
\newcommand{\Subsection}{\subsection}
\providecommand{\nobreakdash}{\nobreak\hbox{-}}
\setlength{\emergencystretch}{2em}
\multlinegap0pt
\usepackage{amssymb}
\usepackage[shortlabels]{enumitem}
\usepackage[normalem]{ulem}
\newtheorem{theorem}{Theorem}
\title{Two remarks on transcendental shift-like maps on $\mathbb{C}^N$}
\author{Ramanpreet Kaur}
\date{Source: arXiv:2604.27832v1 (CC BY 4.0)}
\begin{document}
\maketitle

\noindent\emph{Source and licence.} Two remarks on transcendental shift-like maps on $\mathbb{C}^N$. Version arXiv:2604.27832v1 (CC BY 4.0), distributed by arXiv under the Creative Commons Attribution 4.0 International licence, http://creativecommons.org/licenses/by/4.0/, as recorded in the arXiv licence metadata for this exact version and shown on the abstract page https://arxiv.org/abs/2604.27832v1. Full licence notice: \texttt{licenses/arxiv-2604.27832-CC-BY-4.0.txt}.\par\smallskip

\noindent\textit{Editorial scope.} Counted unit: the article's theorem theorem (source0.tex lines 340--342). The article's complete original proof of it is delivered with it (source0.tex lines 343--370, one \texttt{proof} environment of 28 lines). No mathematics is written, repaired or re-derived: the statement and the proof are the article's own lines, in the article's own notation, and no retained line is substituted or re-flowed.  No further source-local context is needed: every cross-reference of every retained line resolves inside the two delivered spans.  Nothing was omitted from any retained span, so the package carries no omission record.  No retained line contains a citation, so the wrapper carries no bibliography and no bibliography entry.  No retained line contains a figure environment, an image-inclusion command or drawing code, so no image is delivered and the package declares no asset dependency.  Editorial formatting only, in addition: an AMS \texttt{amsart} preamble that restates the article's own declarations and macros used by the retained lines, namely the environments \texttt{theorem} on separate counters, together with the standard packages those lines need.  The retained environments print in this document as Theorem 1 in order of appearance, whereas the article prints the same items with the numbering of its own full document.  The complete native source of the article as distributed in the arXiv e-print for this exact version is bundled unchanged as \texttt{sources/199501/source0.tex}.
\begin{theorem}
The escaping set of a transcendental shift map of type $\nu$ is always non-empty.
\end{theorem}

\begin{proof}
We shall first prove this for $\nu=N-\nu.$ We shall follow Eremenko's approach, that is,......
First observe that if there exists a $z=(z_1,z_2,\dots,z_N)$ such that $\|F^{n\nu}(z)\|\to\infty$ as $n\to\infty$, then $\|F^{n}(z)\|\to\infty$ as $n\to\infty.$ Using Wiman-Valiron estimate, we can choose $R>0$ sufficiently large so that the following properties are satisfied.
 
In this case, note that for $1\leq \nu\leq N-1,$
\[F^{-\nu}(z_1,z_2,\dots,z_N)=\left(\frac{f(x_{N-2\nu+1})-x_{N-\nu+1}}{a},\frac{f(x_{N-2\nu+2})-x_{N-\nu+2}}{a},\dots, \frac{f(x_{N-\nu})-x_{N}}{a}, x_1,x_2,\dots, x_{N-\nu}\right).\]
For every $n\in\mathbb{N},$ define 
\[\Gamma_n=\{(\underbrace{\phi_{n-1}(\phi_n(z)),\phi_n(z),\phi_n(z),\dots,\phi_n(z)}_{(N-v) \text{ times}}, \underbrace{z,z\dots,z}_{\nu \text{ times}}): z\in D_n\}.\]
Observe that for every $n\geq 1$, we have $F^{-\nu}(\Gamma_n)\subset\subset \Gamma_{n-1}$. For this, consider 
\begin{align*}
F^{-\nu}(\phi_{n-1}(\phi_n(z)), \phi_n(z),\phi_n(z),\dots,\phi_n(z), z,z\dots,z)=&\left(\frac{f(\phi_{n}(z))-z}{a}, \frac{f(\phi_{n}(z))-z}{a}, \dots,\frac{f(\phi_n(z))-z}{a},\right.\\
&\left. \phi_{n-1}(\phi_n(z)),z,z,\dots,z\right)\\
=&\left(\phi_{n-1}(\phi_n(z)), \phi_{n-1}(\phi_n(z))\dots,\phi_{n-1}(\phi_n(z)),\right.\\
&\left. \phi_{n-1}(\phi_n(z)), \phi_n(z),\phi_n(z),\dots,\phi_n(z)\right)
\end{align*}
We shall construct an increasing sequence of sets $\{K_n\}$ such that $K_n\subset\subset K_{n+1}$ for every $n\geq 1$.
For this, define 
$\Gamma_{n\nu}=\{(\phi_n(z))\}$\\
Choose $r_0>2$ (sufficiently large) such that it does not belong to the the exceptional set $E$ and satisfies the following properties.
\begin{align}
M(r)&>4r \text{ for every }r\geq r_0\\
|\log(1+\epsilon_i)|&<1  \text{ for every }r\geq r_0 \text{ and }  \text{ for every } i\geq 0.\\
\text{lm}(E\cap [r_0,\infty))&<1\\
N(r_1)&>\max\left\{10^4, \frac{\pi-\frac{\delta}{2}}{10}\right\}
\end{align}
Define
\[S_0=\left\{z\in\mathbb{C}: \left|\log\left|\frac{z}{\zeta_0}\right|\right|<\frac{5}{N(r_0)}, \left|Arg\left|\frac{z}{\zeta_0}\right|\right|<\frac{5}{N(r_0)}\right\}\]
\end{proof}

\end{document}
